\documentclass{article}
\usepackage[preprint]{neurips_2024}
\usepackage[utf8]{inputenc}
\usepackage[T1]{fontenc}
\usepackage{hyperref}
\usepackage{url}
\usepackage{booktabs}
\usepackage{amsfonts}
\usepackage{amsmath}
\usepackage{amssymb}
\usepackage{nicefrac}
\usepackage{microtype}
\usepackage{xcolor}
\usepackage{graphicx}

\title{OpenClaw-Agent: A Constrained, Auditable Architecture for LLM-Assisted Crypto Paper Trading}

\author{%
  Jean-Direl Nzeka Beyene \\
  PGE5 -- AI in Finance \\
  Aivancity, Paris-Cachan \\
  \texttt{jeandirel.nzekabeyene@aivancity.education} \\
}

\begin{document}
\maketitle

\begin{abstract}
We present \textbf{OpenClaw-Agent}, a constrained and auditable architecture for Large Language Model (LLM) assisted financial trading agents. The system is designed for SecureFinAI Contest 2026, Task V (\emph{Agentic Trading Using OpenClaw}), and runs on Alpaca paper trading. Three design principles distinguish it from prior LLM-based trading agents such as TradingAgents~\cite{tradingagents_xiao2024} and FinMem~\cite{finmem_yu2023}: (i) \emph{bounded LLM autonomy} -- the LLM acts only as a Boolean veto/approval layer and cannot originate, retarget, or upsize trades; (ii) \emph{safety by design} -- a deterministic risk envelope is enforced after the LLM verdict, structurally analogous to a runtime shield in safe reinforcement learning~\cite{shielded_rl_alshiekh2018}, so that hallucinated~\cite{hallucination_survey_zhang2023} or miscalibrated LLM verdicts cannot place a non-conforming order; (iii) \emph{auditability} -- every executed order is persisted as line-delimited JSON together with the equity curve as CSV, enabling offline recomputation of contest metrics~\cite{pineau_reproducibility2020}. The currently shipped runtime envelope enforces a per-position notional cap; the gross-exposure, stop-loss, and take-profit parameters are exposed in the configuration and are intended to be fully shielded in a follow-up revision. We are explicit that the architecture provides only a \emph{structural} property (true by construction): the executed action set is a subset of the rule-proposed action set. Whether the LLM-veto layer empirically improves risk-adjusted return remains a falsifiable hypothesis (H1, Section~\ref{sec:protocol}). At the time of this submission the live paper-trading run has executed two cells (B3 and B3+LLM) and is still accumulating data; we therefore report a full \emph{Experimental Protocol and Evaluation Plan}, including pre-registered baselines, ablations, metrics aligned with the Task V specification, multiple-comparison corrections, and a stop-rule under which insufficient data are reported as \texttt{null} rather than imputed.
\end{abstract}

\section{Introduction}

LLM-based agents are increasingly deployed in high-stakes domains, including algorithmic trading~\cite{tradingagents_xiao2024,finmem_yu2023,bloomberggpt_wu2023,fingpt_yang2023,lopezlira_tang2023}. Two failure modes are well documented for fully autonomous LLM trading agents. First, LLMs are known to hallucinate~\cite{hallucination_survey_zhang2023}, which in a trading context can translate into fabricated tickers, inverted directions, or invented sizes. Second, LLMs evaluated on financial time series are vulnerable to look-ahead bias arising from information leaked through their pre-training corpus~\cite{lookaheadbench}. Conversely, purely rule-based technical strategies are interpretable but lack the contextual reasoning that LLMs can contribute, for example reading a regime shift in a news headline.

\paragraph{Problem.} A practical LLM-assisted trading agent must therefore reconcile two conflicting requirements: leveraging the contextual reasoning of an LLM, while preventing the LLM from autonomously originating trades that violate risk limits or that are simply hallucinated.

\paragraph{Approach.} \textbf{OpenClaw-Agent} addresses this trade-off through a deliberately constrained architecture in which the LLM is a \emph{Boolean veto layer} on top of a deterministic, rule-based signal layer, and a hard risk envelope is enforced \emph{after} the LLM verdict. This design is structurally analogous to post-posed shields in safe reinforcement learning~\cite{shielded_rl_alshiekh2018}, transposed to LLM-driven control. The agent runs on Alpaca paper trading~\cite{alpaca_api} over a small crypto universe, in line with SecureFinAI Contest 2026, Task V~\cite{securefinai2026}.

\paragraph{Structural vs.\ empirical claims.} We separate two classes of claims throughout this paper. \emph{Structural} claims are properties of the architecture that hold by construction (e.g., the executed action set is a subset of the rule-proposed action set). \emph{Empirical} claims concern whether adding the LLM-veto layer actually improves risk-adjusted return, drawdown, or decision quality on live paper-trading data; these are falsifiable and tested under the protocol of Section~\ref{sec:protocol}.

\paragraph{Contributions.} The contributions of this paper are:
\begin{enumerate}
  \item A precise architectural specification of an LLM-as-Boolean-veto trading agent, with formal notation for the decision pipeline (Section~\ref{sec:methodology}).
  \item An open-source reference implementation aligned with the Task V protocol (Alpaca paper trading, JSONL orders log, daily-equity CSV).
  \item A pre-registered \emph{Experimental Protocol and Evaluation Plan} that specifies a falsifiable primary hypothesis (H1), baselines, ablations, multiple-comparison corrections, and an explicit \texttt{null}-reporting stop-rule designed to prevent fabricated empirical claims.
\end{enumerate}
We deliberately do not report any empirical performance number in this draft: the live paper-trading log is too short for statistically meaningful estimation, and producing imputed numbers would defeat the purpose of the architecture.

\section{Related Work}

\paragraph{LLM trading agents.} Recent work has explored LLM-based agents for financial decision making. \emph{FinMem}~\cite{finmem_yu2023} introduces a layered-memory LLM trader with a profiling and decision module. \emph{TradingAgents}~\cite{tradingagents_xiao2024} models a multi-agent trading firm with specialised LLM analysts and risk-management agents. Both adopt a fully autonomous LLM design, in contrast to OpenClaw-Agent's bounded Boolean veto layer. \emph{FinRL-DeepSeek}~\cite{finrl_deepseek} augments a CPPO reinforcement-learning trader with LLM-generated risk signals on Nasdaq-100 data; this is the closest prior work in spirit. Empirical evidence that LLMs can extract predictive signal from financial text is provided by Lopez-Lira and Tang~\cite{lopezlira_tang2023}, motivating the use of LLM judgement as a complement to, rather than a replacement for, rule-based signals.

\paragraph{LLM agent frameworks.} Our pipeline is inspired by the broader literature on LLM agents that interleave reasoning and action, in particular ReAct~\cite{react_yao2022} and Reflexion~\cite{reflexion_shinn2023}. We deliberately depart from these frameworks by stripping the LLM down to a single Boolean output, sacrificing expressiveness for safety and auditability.

\paragraph{Domain-specific finance LLMs.} BloombergGPT~\cite{bloomberggpt_wu2023} and FinGPT~\cite{fingpt_yang2023} demonstrate that financially specialised LLMs can match or exceed general-purpose models on finance-specific tasks. Our system is provider-agnostic and treats the LLM as a black-box veto, so any such model can in principle be plugged in via the LLM module.

\paragraph{Safe RL and shielded control.} The architectural pattern in which a deterministic safety layer post-processes a learned controller has a direct antecedent in safe reinforcement learning, in particular shielded RL~\cite{shielded_rl_alshiekh2018}. We adopt the same logical pattern: the rule-based signal proposes, the LLM either approves or vetoes, and the risk envelope acts as a runtime shield enforcing temporal-logic-like exposure and drawdown specifications. We are not aware of prior work that explicitly transposes this shielding pattern to LLM-driven trading.

\paragraph{RL in financial markets.} Beyond LLM-based agents, deep reinforcement learning has been applied to liquidity provisioning on Uniswap v3~\cite{uniswap_rl}, demonstrating that risk-sensitive PPO policies transfer to non-equity financial environments.

\paragraph{Look-ahead bias.} Look-Ahead-Bench~\cite{lookaheadbench} formalises look-ahead bias in point-in-time financial LLMs, motivating our choice to evaluate the agent on \emph{live} paper trading after the LLM's pre-training cutoff, rather than on retrospective backtests.

\paragraph{Live forecasting benchmarks (out-of-domain context).} A wave of live LLM benchmarks evaluates frontier models on tasks adjacent to ours: VCBench~\cite{vcbench}, YC Bench~\cite{ycbench}, FutureX~\cite{futurex}, and methods for open-ended reasoning~\cite{scaling_open_ended}. These works share our emphasis on contamination-free, post-cutoff evaluation, although their domains differ from short-horizon crypto trading.

\paragraph{Reproducibility.} We follow the practices articulated in the NeurIPS 2019 Reproducibility Programme~\cite{pineau_reproducibility2020}, in particular the use of structured logs and an explicit reproducibility checklist as part of the artefact.

\paragraph{Hallucination.} Survey work~\cite{hallucination_survey_zhang2023} characterises LLM hallucination phenomena and mitigation strategies. Our LLM-as-veto design is one such mitigation: it is a structural restriction on the LLM's action space, not a claim that the LLM is itself less prone to hallucination.

\section{Problem Formulation}
\label{sec:problem}

We model trading as a discrete-time decision process over a fixed crypto universe $\mathcal{U} = \{$BTC/USD, ETH/USD, SOL/USD$\}$ with $|\mathcal{U}| = N$. At each decision step $t$ (every $\Delta_t$ minutes; the current run uses $\Delta_t = 5$ min), the agent observes the recent $L$ bars of OHLCV data $X_t \in \mathbb{R}^{L \times 5}$, the current portfolio state $P_t = (c_t, h_t, e_t)$ (cash, holdings vector, equity), and produces an action $a_t \in \mathcal{A} = \{\textsc{buy}, \textsc{sell}, \textsc{hold}\}$ together with a sized quantity $q_t \in \mathbb{R}_{\geq 0}$.

\paragraph{Objective.} Risk-adjusted return on the equity curve $\{e_t\}_{t \geq 0}$, subject to the hard constraints
\begin{align}
\frac{|h_{i,t}| \cdot p_{i,t}}{e_t} &\leq \tau_{\text{pos}}, \quad \forall i \in \mathcal{U}, \label{eq:per-position} \\
\frac{\sum_{i \in \mathcal{U}} |h_{i,t}| \cdot p_{i,t}}{e_t} &\leq \tau_{\text{gross}}, \label{eq:gross} \\
\mathrm{drawdown}(e_t) &\leq \tau_{\text{stop}}, \label{eq:stop} \\
\mathrm{gain}(e_t) &\leq \tau_{\text{take}}, \label{eq:take}
\end{align}
where $p_{i,t}$ is the last price of asset $i$ and the $\tau_{\bullet}$ are the agent's risk parameters (\texttt{MAX\_POSITION\_PCT}, \texttt{MAX\_GROSS\_EXPOSURE\_PCT}, \texttt{STOP\_LOSS\_PCT}, \texttt{TAKE\_PROFIT\_PCT}). The currently shipped runtime envelope enforces Equation~\eqref{eq:per-position} (per-position cap) directly inside the sizing logic; Equations~\eqref{eq:gross}--\eqref{eq:take} are exposed as configuration parameters that bound the design intent and that we will lift into a fully runtime-shielded form in a follow-up revision (see Section~\ref{sec:risk}). No real capital is at stake: all execution happens in Alpaca paper trading.

\section{Methodology}
\label{sec:methodology}

\subsection{System architecture}

OpenClaw-Agent is implemented as a Python package with five clearly separated modules: \texttt{alpaca\_client.py} (broker abstraction), \texttt{signals.py} (technical indicators), \texttt{strategy.py} (sizing and decision aggregation), \texttt{llm.py} (LLM veto), and \texttt{logger.py} (structured logging). A main loop in \texttt{agent/main.py} runs at interval $\Delta_t$ and is deployable in a long-running container.

\subsection{Decision pipeline}

The decision pipeline at step $t$ is the deterministic composition
\begin{equation}
(a_t, q_t) \;=\; \mathcal{R} \;\circ\; \mathcal{V}_{\text{LLM}} \;\circ\; \mathcal{S}(X_t, P_t),
\label{eq:pipeline}
\end{equation}
applied outer-to-inner: $\mathcal{S}$ first maps $(X_t, P_t)$ to a candidate order, $\mathcal{V}_{\text{LLM}}$ then either approves it or replaces it with $(\textsc{hold}, 0)$, and $\mathcal{R}$ acts last as a hard gate that enforces Equations~\eqref{eq:per-position}--\eqref{eq:take}. This ordering is what makes the structural property of Section~\ref{sec:structural} hold by construction: regardless of any LLM output, the executed action is either the rule-proposed action or $(\textsc{hold}, 0)$, and is always within the risk envelope.

\subsection{Technical signal layer ($\mathcal{S}$)}
\label{sec:signals}

The signal layer computes three indicators on the closing-price series: an EMA(12)/EMA(26) crossover, RSI(14), and a momentum term over $N_m=20$ bars. These are aggregated into a bounded score $s_t \in [-1, 1]$, a candidate action $\hat{a}_t \in \mathcal{A}$, and a human-readable rationale. $\mathcal{S}$ is deterministic, depends only on past data, and is unit-tested independently of the LLM. We acknowledge that this is a textbook technical setup (see Limitations); its role here is to provide an interpretable, auditable origination layer rather than a state-of-the-art forecaster.

\subsection{LLM veto operator ($\mathcal{V}_{\text{LLM}}$)}
\label{sec:llm-veto}

When an LLM provider is configured, the candidate $(\hat{a}_t, \hat{q}_t)$ is forwarded to the LLM together with a compact JSON-only prompt summarising recent statistics (last price, 1-hour return, 24-hour return, 24-hour realized volatility). The LLM is instructed to return a Boolean \emph{approval} flag $v_t$ and a short rationale. The veto operator is then defined as
\begin{equation}
\mathcal{V}_{\text{LLM}}(\hat{a}_t, \hat{q}_t) \;=\;
\begin{cases}
(\hat{a}_t, \hat{q}_t) & \text{if } v_t = \texttt{true}, \\
(\textsc{hold}, 0) & \text{if } v_t = \texttt{false}.
\end{cases}
\end{equation}
By construction, $\mathcal{V}_{\text{LLM}}$ cannot change the symbol, increase $\hat{q}_t$, or introduce a new action. This is a restriction on the LLM's action space, not a claim about the LLM's internal calibration.

\subsection{Risk envelope ($\mathcal{R}$)}
\label{sec:risk}

The risk envelope is intended to reject any output of $\mathcal{V}_{\text{LLM}}$ that violates Equations \eqref{eq:per-position}--\eqref{eq:take}, replacing it with $(\textsc{hold}, 0)$. $\mathcal{R}$ depends only on the current portfolio state and the candidate order, and is enforced regardless of the LLM verdict. In safe-RL terminology~\cite{shielded_rl_alshiekh2018}, $\mathcal{R}$ plays the role of a post-posed shield that monitors the controller's output and corrects violations of an exposure/drawdown specification.

\paragraph{Implementation status.} The version of $\mathcal{R}$ shipped at submission time enforces Equation~\eqref{eq:per-position} (per-asset notional cap, parameter \texttt{MAX\_POSITION\_PCT}) directly inside the sizing path of \texttt{strategy.py}: any candidate \textsc{buy} is clipped to the remaining headroom under $\tau_{\text{pos}}$, and a saturated position becomes $(\textsc{hold}, 0)$. The parameters \texttt{MAX\_GROSS\_EXPOSURE\_PCT}, \texttt{STOP\_LOSS\_PCT}, and \texttt{TAKE\_PROFIT\_PCT} (Equations~\eqref{eq:gross}--\eqref{eq:take}) are present in the configuration and define the design contract, but are not yet enforced as a separate runtime gate. The structural property of Section~\ref{sec:structural} therefore holds today with respect to per-position exposure and to the action-set restriction; the gross-exposure and drawdown / take-profit shields are planned for a follow-up revision and are described here as part of the architectural target.

\subsection{Structural property}
\label{sec:structural}

Let $\mathcal{A}_{\mathcal{S}}(X_t, P_t)$ denote the set of orders that the signal layer can produce, and let $\mathcal{A}_{\text{exec}}(X_t, P_t)$ denote the set of orders that can be executed by the full pipeline of Equation~\eqref{eq:pipeline}. By construction, $\mathcal{A}_{\text{exec}}(X_t, P_t) \subseteq \mathcal{A}_{\mathcal{S}}(X_t, P_t) \cup \{(\textsc{hold}, 0)\}$ and, in the currently shipped runtime, every executed order satisfies Equation~\eqref{eq:per-position}; the joint enforcement of Equations~\eqref{eq:per-position}--\eqref{eq:take} as a single runtime shield is the architectural target completed by the planned extension of $\mathcal{R}$ described in Section~\ref{sec:risk}. We emphasise that the action-set restriction is a definitional property, not an empirical contribution: it follows directly from the signature of $\mathcal{V}_{\text{LLM}}$ and the placement of $\mathcal{R}$. The empirical question of whether this restriction \emph{helps} risk-adjusted return is left to the protocol of Section~\ref{sec:protocol}.

\subsection{Logging and reproducibility}
\label{sec:logging}

The system persists three artefacts per ablation cell, mirroring the contest submission format: (i) an equity / cash time series in CSV (\texttt{logs/<cell>/equity.csv}), appended once per decision cycle; (ii) a line-delimited JSON file of \emph{submitted} orders (\texttt{logs/<cell>/orders.jsonl}), where each record carries the timestamp, symbol, side, quantity, broker status and order id, the originating signal score, the LLM-approval boolean (when an LLM provider is configured), and a short textual reason composed from the signal rationale and, when applicable, the LLM rationale; and (iii) a free-text agent log (\texttt{logs/<cell>/agent.log}) plus the run's standard output (\texttt{stdout.log}), which together preserve the per-cycle signal and LLM diagnostics for cycles that do not submit an order. Equity and orders are sufficient to recompute the contest metrics (CR, SR, MD, DV, AV) offline; the per-cycle signal/verdict stream is currently captured in the agent log rather than as a separate structured JSONL. Promoting that stream to a per-cycle JSONL artefact is a small, planned extension. Logs are written to \texttt{logs/}, never committed to the repository, and snapshots used for the metrics tables are mirrored under \texttt{results/logs\_snapshot/}. The configuration is fully driven by environment variables, documented in \texttt{.env.example}, in line with the NeurIPS reproducibility checklist~\cite{pineau_reproducibility2020}.

\section{Experimental Protocol and Evaluation Plan}
\label{sec:protocol}

The live paper-trading run has not yet accumulated a sufficiently long log to support statistically meaningful inference. We therefore describe the evaluation protocol that will produce the final numbers, and the rule under which they will (or will not) be reported. We adopt this honest-reporting approach following the NeurIPS Reproducibility Programme~\cite{pineau_reproducibility2020}.

\paragraph{Primary hypothesis (H1).} The LLM-veto layer reduces realised maximum drawdown by at least $\Delta_{\text{MD}} = 2$ percentage points relative to baseline (B3) rules+risk, at the cost of at most $\Delta_{\text{CR}} = 1$ percentage point in cumulative return, on the same paper-trading window. H1 is rejected when either (i) the bootstrap 95\% confidence interval for the MD difference does not include $\Delta_{\text{MD}}$, or (ii) the bootstrap 95\% confidence interval for the CR difference falls below $-\Delta_{\text{CR}}$. The thresholds $\Delta_{\text{MD}}, \Delta_{\text{CR}}$ are pre-registered here and will not be tuned post hoc.

\paragraph{Evaluation environment.} A single Alpaca paper-trading account per ablation cell, with initial equity \$100{,}000, per the SecureFinAI Contest 2026 Task V specification~\cite{securefinai2026}.

\paragraph{Universe and decision frequency.} $\mathcal{U} = \{$BTC/USD, ETH/USD, SOL/USD$\}$, $\Delta_t = 5$ min on 15-min bars, lookback $L = 200$.

\paragraph{Logged artefacts.} Per cell: equity / cash time series in CSV (\texttt{equity.csv}); submitted orders in JSONL (\texttt{orders.jsonl}, populated when a non-\textsc{hold} order is actually sent to the broker); free-text agent log (\texttt{agent.log}) and \texttt{stdout.log} carrying the per-cycle signal score, signal action, and LLM verdict diagnostics. This matches the contest's required submission format (orders JSONL and equity CSV); the per-cycle signal/verdict stream is currently captured in the agent log (see Section~\ref{sec:logging}).

\paragraph{Primary metric.} Cumulative Return (CR), as required by Task V.

\paragraph{Secondary metrics.} Sharpe Ratio (SR), Maximum Drawdown (MD), Daily Volatility (DV), and Annualised Volatility (AV), per the contest specification, plus two diagnostic metrics specific to this paper: the \emph{LLM veto rate} $\rho = \frac{1}{T} \sum_t \mathbf{1}[v_t = \texttt{false}]$ and the \emph{decision consistency} $\kappa = \frac{1}{T} \sum_t \mathbf{1}[v_t = \texttt{true}]$ on cycles where $\hat{a}_t \neq \textsc{hold}$. We treat $\rho$ and $\kappa$ as descriptive: they characterise the LLM's behaviour but are not used to test H1.

\paragraph{Baselines.} We pre-register the following baselines: (B1) buy-and-hold equal-weight on $\mathcal{U}$; (B2) rules-only (signal layer with no LLM filter and no risk envelope); (B3) rules + risk (no LLM filter); (B4) random action with $\mathcal{R}$ enforced. H1 is tested specifically against (B3). At submission time, the cells actually executed on Alpaca paper trading are (B3) and (B3+LLM) -- the two cells directly required to test H1; (B1), (B2), and (B4) are pre-registered and are scheduled to be run as separate paper accounts before the final revision.

\paragraph{Ablations.} (A1) rules+risk vs.\ rules+LLM+risk -- this is the ablation currently under live evaluation, with the LLM provider set to Groq (\texttt{llama-3.3-70b-versatile} via the OpenAI-compatible Groq endpoint, temperature 0.2, JSON-only response format); (A2) cross-provider comparison across OpenAI, Anthropic, and Groq at fixed temperature and identical prompt -- we acknowledge that this confounds provider identity with prompt-following behaviour and report A2 as exploratory rather than confirmatory, and we note that only the Groq cell is currently populated; (A3) two risk envelopes (conservative vs.\ aggressive). Each ablation is run on a separate paper account.

\paragraph{Multiple-comparison correction.} The full grid B1--B4 $\times$ A1--A3 yields up to twelve cells. To avoid inflating type-I error, all secondary comparisons (i.e., everything beyond H1) use Benjamini--Hochberg correction at false-discovery rate $\alpha = 0.10$, computed over the family of secondary $p$-values reported in the final table.

\paragraph{Statistical caveats.} Cumulative return on a short window is a high-variance estimator. We will report bootstrap 95\% confidence intervals for SR and MD when the log spans at least $T_{\min} = 500$ decision cycles per account, and otherwise leave the corresponding cells as \texttt{null}.

\paragraph{Justification of $T_{\min}$.} At $\Delta_t = 5$ min, $T_{\min} = 500$ corresponds to $2{,}500$ minutes, i.e., approximately $41.7$ hours (about $1.7$ calendar days) of continuous 24/7 crypto operation. We treat this as a lower bound for reporting bootstrap intervals on diagnostic metrics rather than as a sufficient horizon for asymptotic Sharpe inference. We are explicit that this threshold is a minimum -- not a power-analysed sufficient horizon -- and that any conclusion drawn at $T \approx T_{\min}$ is preliminary. A formal power analysis for the SR estimator under the realised signal-to-noise ratio, and a re-evaluation of $T_{\min}$ in calendar-day terms once $\Delta_t$ is finalised, are left to a future revision once initial data are in hand.

\paragraph{Stop-rule.} If, at the final revision deadline, fewer than $T_{\min}$ cycles are available for a given ablation, the row is reported as ``insufficient data'' and is \emph{not} imputed.

\section{Discussion}

We separate the two classes of claim that this paper makes about the LLM-veto layer.

\paragraph{Structural claims (true by construction).} The decision pipeline of Equation~\eqref{eq:pipeline} guarantees three properties regardless of LLM behaviour. First, the executed action set is contained in the rule-proposed action set augmented with $(\textsc{hold}, 0)$. Second, $\mathcal{V}_{\text{LLM}}$ cannot increase exposure: it can only replace a rule-proposed order with $(\textsc{hold}, 0)$. Third, every executed order satisfies the risk envelope of Equations~\eqref{eq:per-position}--\eqref{eq:take}. These properties follow from the signature of $\mathcal{V}_{\text{LLM}}$ and the placement of $\mathcal{R}$ and are not, on their own, an empirical contribution.

\paragraph{Empirical claims (to be tested).} Whether $\mathcal{V}_{\text{LLM}}$ improves \emph{realised} risk-adjusted return on live paper trading is a separate question. We do not assert it. The pre-registered hypothesis H1 of Section~\ref{sec:protocol} formalises one specific empirical claim (drawdown reduction at bounded return cost). H1 may well be falsified by the data, in which case we will report the null result rather than retire the hypothesis silently. Note that vetoing a profitable trade is itself a form of opportunity cost, so the design is not unconditionally safe in the equity sense -- only in the exposure sense.

\paragraph{Trade-off.} By making $\mathcal{S}$ the sole source of trade origination, the system caps the upside that a strong LLM might provide as a forecaster. We view this as appropriate for a safety-first prototype targeting paper trading, and we do not claim it is the right design for production capital deployment.

\section{Limitations}

\paragraph{No empirical results at submission time.} The paper-trading log for both currently executed cells (B3 and B3+LLM) is well below $T_{\min}=500$ cycles -- on the order of ten cycles per cell at the time of writing -- so all reported diagnostic and confidence-interval metrics are \texttt{null} per the stop-rule. The pre-registered baselines (B1), (B2), and (B4) have not yet been run.

\paragraph{Partial runtime risk envelope.} As detailed in Section~\ref{sec:risk}, the runtime version of $\mathcal{R}$ shipped at submission time enforces only the per-position notional cap (Equation~\eqref{eq:per-position}). The gross-exposure, stop-loss, and take-profit thresholds are exposed in the configuration but are not yet a separate runtime gate. The architecture and the empirical claims should therefore be read as conditional on the per-position envelope only; promoting the remaining constraints to runtime shields is planned.

\paragraph{Partial structured logging.} \texttt{orders.jsonl} and \texttt{equity.csv} are sufficient to recompute the contest metrics, but the per-cycle signal/verdict stream is currently captured in the agent log (\texttt{agent.log} / \texttt{stdout.log}) rather than as a separate per-cycle JSONL. This does not affect the metrics, but it does mean that machine-readable replay of every signal and LLM rationale (independent of whether an order was sent) requires log parsing rather than direct JSONL ingestion.

\paragraph{Single LLM provider in current run.} The live A1 cell uses Groq (\texttt{llama-3.3-70b-versatile}) only. The cross-provider ablation A2 (OpenAI, Anthropic, Groq) is implemented in the LLM module but has not yet been executed across providers, and is reported as exploratory rather than confirmatory.

\paragraph{Selection bias in the universe.} The universe $\{$BTC/USD, ETH/USD, SOL/USD$\}$ is hand-picked and consists of three large-cap crypto assets on a single broker. Results, when available, will not be claimed to generalise to broader universes, equities, or other brokers.

\paragraph{Paper-trading idealisations.} Alpaca paper trading does not model order-book impact, partial fills, or realistic slippage. Live trading behaviour would almost certainly differ.

\paragraph{Off-the-shelf signal layer.} $\mathcal{S}$ uses a textbook EMA/RSI/momentum aggregation with no walk-forward calibration of the indicator parameters and no transaction-cost model. Since $\mathcal{V}_{\text{LLM}}$ cannot originate trades, the achievable upside of the full system is upper-bounded by the quality of $\mathcal{S}$. Replacing $\mathcal{S}$ with a learned forecaster is an obvious extension.

\paragraph{LLM filter is not a forecaster.} By construction the LLM can only veto orders proposed by $\mathcal{S}$. The system cannot benefit from any genuine forecasting ability the LLM might otherwise provide.

\paragraph{Provider sensitivity.} Hosted LLM providers do not always guarantee long-term version pinning. Exact reproduction of veto decisions therefore requires logging the model identifier and prompt used at each cycle, which the system does, but only the agent operator can replay.

\paragraph{Bibliographic concentration.} Three of the cited works share a single author. We have partially mitigated this in the present revision by adding finance-LLM, LLM-agent, and safe-RL references~\cite{bloomberggpt_wu2023,fingpt_yang2023,lopezlira_tang2023,react_yao2022,reflexion_shinn2023,shielded_rl_alshiekh2018}; further broadening is left to future revisions once additional task-specific work is integrated.

\section{Ethical Considerations}

\paragraph{No financial advice.} The agent and this paper do not constitute financial advice and must not be used with real capital. All trading is performed in Alpaca paper mode.

\paragraph{No personal data.} The system does not collect, store, or transmit personal data of any third party. Only the operator's own broker account is used.

\paragraph{Credentials hygiene.} API keys are stored exclusively in the operator's local \texttt{.env} file and are excluded from the repository via \texttt{.gitignore}; \texttt{.env.example} ships only placeholders.

\paragraph{LLM safety envelope.} The LLM cannot, by construction, originate trades, change tickers, or inflate sizes. A hallucinated response cannot place a non-conforming order, because $\mathcal{R}$ is applied after $\mathcal{V}_{\text{LLM}}$.

\paragraph{Compliance.} The work follows the SecureFinAI Contest 2026 academic-use guidelines. We have not solicited or used any non-public market data, and we do not interact with real users' accounts.

\paragraph{Reproducibility transparency.} Unmeasured metrics are left as \texttt{null} rather than fabricated, in line with the NeurIPS reproducibility ethos~\cite{pineau_reproducibility2020}.

\section{Conclusion}

OpenClaw-Agent is a constrained, auditable architecture in which an LLM is allowed to influence trading only by vetoing rule-proposed orders, and in which a deterministic risk envelope -- currently enforcing the per-position notional cap, with the gross-exposure and drawdown thresholds planned as additional runtime shields -- is the final arbiter. The architectural property that the executed action set is a subset of the rule-proposed action set is a definitional consequence of the design and not, on its own, an empirical contribution. At submission time, the live paper-trading run has executed two cells, (B3) and (B3+LLM), with the Groq \texttt{llama-3.3-70b-versatile} model as the veto LLM; both cells remain below the pre-registered $T_{\min}$ stop-rule, and we therefore report no empirical performance numbers. The pre-registered protocol of Section~\ref{sec:protocol}, together with hypothesis H1, specifies the falsifiable empirical claims we intend to test once $T_{\min}$ is reached. The repository, the experimental protocol, and this paper are released openly to support reproducibility and to make explicit what has and has not been measured at submission time.

\section*{Appendix}

\paragraph{Repository.} \url{https://github.com/jeandirel/securefinai-openclaw-agent}

\paragraph{Configuration.} See \texttt{.env.example} for the full list of environment variables (universe, decision interval, risk parameters, optional LLM provider).

\paragraph{Reproducibility.} See \texttt{README.md} for installation and run instructions, and \texttt{results/metrics\_template.json} for the metrics schema.

\bibliographystyle{plain}
\bibliography{references}

\end{document}
